% Working section of the combined HCP and proof-carrying transaction manuscript.
% Canonical prose counterpart: ../DRAFT.md, Sections 2.1 and 6.3.1.
% Requires amsmath, amssymb, amsthm and proposition/theorem environments.
\section{Coupled harmonic lattices and their public class}
\label{sec:coupled-harmonic-construction}

Homomorphic correctness requires evaluation on ciphertexts to agree with
evaluation on messages \cite{gentry2009fullyHomomorphic}. The homotopical
analogue developed here concerns admissible presentation transport.
For an integral chain isometry compatible with the harmonic-lattice encodings,
the existing sample-level transport and recovery theorems give
\[
 T_1\operatorname{Enc}_F(z,e)
   =\operatorname{Enc}_G(z,T_{\mathrm{noise}}e),\qquad
 \operatorname{Dec}_G(T_1b)=z=\operatorname{Dec}_F(b)
\]
on encryptions of \(z\). The underlying cochain equivalence preserves the
declared semantic homotopy class. This naturality square states the proved
analogy. Transport material may be private; the equation alone supplies no
public translation algorithm. Ciphertext transformation under changed
public keys also appears in fully key-homomorphic encryption
\cite{bonehEtAl2014keyHomomorphic}. Our comparison concerns the concrete
presentation data, permitted transports and their semantic class.

The cellular sheaf determines the cochain complex before a message is chosen.
For an integral presentation, the admissible messages form the subgroup
\[
  \Lambda_F=C^1_{\mathbb Z}(F)\cap\ker\Delta_{1,F}\ \subseteq\ C^1(F).
\]
A message in this lattice is already a degree-one cochain. An exact/coexact
opening produces the observed cochain
\begin{equation}
  b=h+\delta_0a+\delta_1^*u,\qquad h\in\Lambda_F,
  \label{eq:coupled-hodge-encryption}
\end{equation}
and the original harmonic projector recovers \(h\). The lattice specifies the
integral message space; preservation of the public homotopy class comes from
an equivalence of the underlying cochain interpretations.

\subsection{Shared integral stalk constraints}

Fix \(B\in\mathbb Z^{m\times n}\). The coupled cellular family has one scalar
exact edge, \(n\) scalar information edges, \(m\) scalar parity edges, and
one scalar face stalk per row of \(B\). Its actual restriction maps give
\[
 C^1(F_B)=\mathbb R\oplus\mathbb R^n\oplus\mathbb R^m,
 \qquad
 \delta_{0,B}(a)=(a,0,0),\qquad
 \delta_{1,B}(a,x,y)=y-Bx.
\]
The information edges contribute to shared face equations, so the stalk maps
couple different edge cells. Standard integer stalk lattices are preserved by
these maps. All cochain spaces use their fixed Euclidean metrics.

\begin{proposition}[Intrinsic harmonic graph lattice]
\label{prop:coupled-harmonic-graph}
The harmonic subspace and its intrinsic integral lattice are exactly
\[
 W_B=\{(0,x,Bx):x\in\mathbb R^n\},\qquad
 \Lambda_B=\{(0,z,Bz):z\in\mathbb Z^n\}.
\]
The columns of \(G_B=[0;I;B]\) present the entire lattice \(\Lambda_B\), with
Gram matrix \(G_B^{\mathsf T}G_B=I+B^{\mathsf T}B\). For target \((a,v,w)\),
the squared distance to the point with coefficient vector \(z\) is
\[
 a^2+\|v-z\|^2+\|w-Bz\|^2.
\]
The embedded harmonic subspace determines \(B\).
\end{proposition}

\begin{proof}
The harmonic equations are \(\delta_{0,B}^*(a,x,y)=0\) and
\(\delta_{1,B}(a,x,y)=0\), which give \(a=0\) and \(y=Bx\).
Integrality of the information coordinates forces \(x\in\mathbb Z^n\);
integer \(B\) then makes every parity coordinate integral. The information
coordinates also prove independence of the graph generators. Their inner
products and distances give the displayed formulas. Applying the graph
description to the information-coordinate unit vectors recovers each column
of \(B\).
\end{proof}

Real graph matrices describe a standard information-coordinate chart of a
Grassmannian; integer matrices select arithmetic points in that chart. The
formalization constructs the sheaves, harmonic subspaces and lattices, without
an algebraic-variety or scheme construction. Shared rows give explicit
off-diagonal Gram entries, while computational hardness requires an
independent reduction and a specified experiment.

\subsection{Constructed homotopy-class and lattice preservation}

\begin{theorem}[Same public class and message coefficients]
\label{thm:coupled-public-class}
For integer matrices \(B,C\) of the same dimensions, define
\[
 T_{B,C}(a,x,y)=(a,x,y+(C-B)x).
\]
Together with identity maps in degrees zero and two, this is a strict
isomorphism between the actual cellular cochain complexes. It induces a
typed Helfer equivalence and equality of their closed theories in the declared
homological test fragment. Its restriction gives an additive equivalence
\(\Lambda_B\simeq\Lambda_C\) preserving the integer coefficient vector:
\[
 T_{B,C}(0,z,Bz)=(0,z,Cz).
\]
\end{theorem}

\begin{proof}
The inverse is \(T_{C,B}\), and direct substitution gives
\[
 T_{B,C}\delta_{0,B}=\delta_{0,C},\qquad
 \delta_{1,C}T_{B,C}=\delta_{1,B}.
\]
The resulting strict complex isomorphism supplies the typed semantic
equivalence. Proved invariance of the declared closed tests gives theory
equality. The integral graph description in
Proposition~\ref{prop:coupled-harmonic-graph} gives the lattice equivalence
and its coefficient formula.
\end{proof}

This is an isomorphism of the full cochain interpretations. The shear mixes
edge cells and need not define a stalkwise cellular morphism on the fixed
scalar-cell decomposition. It generally changes lengths. Consequently it
does not supply naturality of adjoints or orthogonal projections, or transport
the noise in~\eqref{eq:coupled-hodge-encryption}. The protocol's
\texttt{PresentedRotation} retains its stronger integral-isometry and
same-coefficient conditions. The admitted semantic family can therefore be
larger than a proved protocol rotation orbit.

The realized HCP matrix acts by the actual cellular Hodge Laplacian. Its
matrix kernel equals \(W_B\), and its operative recovery map equals the
original cellular harmonic projector. Thus semantic interpretation, harmonic
lattice, message and observed cochain refer to the same construction. Equal
public classes leave equality of released roots, headers and proof transcripts
as a separate application obligation.

\subsection{Canonical presentation bytes and storage}

The versioned presentation domain \texttt{HODGE\_COUPLING\_V1} fixes this
family's scalar stalks, standard integral structures and Euclidean metric.
The bytes record both dimensions followed by the signed integer entries in
row-major order. Natural magnitudes use little-endian base-256 digits, each
tagged to distinguish it from the terminating byte. An integer adds a sign
byte. Canonical re-encoding rejects redundant digits and negative zero.
The executable decoder roundtrips the actual matrix; injectivity, rejection
of wrong dimensions or domain, trailing data and every proper truncation are
proved. Encoded magnitude length is logarithmic in its value.

An optional storage profile equates the canonical presentation hash with the
original sheaf-indexed snapshot root. A successful parser and root check
recovers the named matrix packet or exhibits distinct byte strings with equal
hashes. Snapshot authentication, payload-envelope binding and agreement with
the semantic classifier retain their original explicit contracts. The
presentation bytes are private witness material in the hidden-operator
profile. Releasing those bytes reveals \(B\); a digest alone establishes no
hiding property.

The exact payload encodes \(\delta_0a\), whereas the full native HCP ciphertext
is \(b\) in~\eqref{eq:coupled-hodge-encryption}. The two objects retain distinct
bindings. A proximity theorem about the exact payload does not automatically
observe the harmonic graph.

For the constructed integral opening profile, \(z,a,u\) are integer-valued and
the original Hodge ciphertext has coordinates
\[
 b=(a,\ z-B^{\mathsf T}u,\ Bz+u).
\]
An executable integer vector casts to exactly this original observation. The
full-ciphertext codec uses the distinct domain
\texttt{HODGE\_CIPHERTEXT\_V1}, explicit matrix dimensions, and exact,
information and parity coordinates in that order. Its decoder reconstructs
the same native sample's full observed cochain; the existing harmonic
projector recovers its original message. Canonical roundtrip, injectivity and
rejection of malformed or mismatched byte streams are proved. A separate
ciphertext hash check binds these bytes to the existing bound observation,
retaining its exact artifact, snapshot and lifecycle. A concrete proof of the
selected transaction/proximity relation remains an application obligation.

The constructed operations concern harmonic-lattice messages, their Hodge
openings and recovery. Encrypted linear query evaluation is an optional
application. A later comparison with fully homomorphic encryption should
inventory the actual supported operations and their security statements.

The shared parity register records the construction, encodings, transport
square and their formal boundaries as H81--H88.
